\documentclass[12pt,a4paper]{article}

\usepackage[utf8]{inputenc}
\usepackage[spanish]{babel}
\usepackage{geometry}
\geometry{top=2.5cm, bottom=2.5cm, left=2.5cm, right=2.5cm}
\usepackage{bookmark}
\usepackage{hyperref}
\usepackage{xcolor}
\usepackage{listings}
\usepackage{graphicx}
\usepackage{titlesec}
\usepackage{tcolorbox}
\usepackage{parskip}
\usepackage{booktabs}
\usepackage{tabularx}

% Colors
\definecolor{primary}{RGB}{41, 128, 185}
\definecolor{secondary}{RGB}{44, 62, 80}
\definecolor{accent}{RGB}{231, 76, 60}
\definecolor{codebg}{RGB}{245, 245, 245}

% Setup links
\hypersetup{
    colorlinks=true,
    linkcolor=primary,
    filecolor=magenta,      
    urlcolor=primary,
    pdftitle={Documentación Técnica: Backend Mecánico},
}

% Code snippets config
\lstdefinestyle{mystyle}{
    backgroundcolor=\color{codebg},   
    commentstyle=\color{green!50!black},
    keywordstyle=\color{blue},
    numberstyle=\tiny\color{gray},
    stringstyle=\color{orange},
    basicstyle=\ttfamily\footnotesize,
    breakatwhitespace=false,         
    breaklines=true,                 
    captionpos=b,                    
    keepspaces=true,                 
    numbers=left,                    
    numbersep=5pt,                  
    showspaces=false,                
    showstringspaces=false,
    showtabs=false,                  
    tabsize=2
}
\lstset{style=mystyle}

% Titles styling
\titleformat{\section}{\color{secondary}\Large\bfseries}{\thesection}{1em}{}[\vspace{-0.5ex}\rule{\textwidth}{1pt}]
\titleformat{\subsection}{\color{primary}\large\bfseries}{\thesubsection}{1em}{}

\begin{document}

\begin{titlepage}
    \centering
    \vspace*{2cm}
    
    {\Huge \textbf{\textcolor{secondary}{Apex Forge}}}\\[0.5cm]
    {\LARGE \textbf{Documentación Técnica de la API Backend}}\\[1.5cm]
    
    {\Large \textit{Sistema de Gestión y Agendamiento para Taller Mecánico}}\\[2cm]
    
    \begin{tcolorbox}[colback=primary!10, colframe=primary, arc=5mm, width=0.8\textwidth, center]
        \centering
        \vspace{0.5cm}
        \textbf{Stack Tecnológico:}\\
        Node.js $\cdot$ NestJS $\cdot$ TypeScript $\cdot$ Prisma ORM\\
        PostgreSQL $\cdot$ Docker $\cdot$ JWT $\cdot$ Render
        \vspace{0.5cm}
    \end{tcolorbox}
    
    \vfill
    {\large \today}\\[1cm]
\end{titlepage}

\tableofcontents
\newpage

\section{Resumen Ejecutivo}
Este documento describe el desarrollo y la arquitectura del backend para el sistema de agendamiento de \textbf{Apex Forge}. El objetivo principal del sistema es proporcionar una plataforma robusta orientada a servicios (SOA) que permita a los clientes agendar citas automotrices, mientras dota a los administradores/mecánicos de una interfaz de gestión para aprobar citas, modificar estados y bloquear días en el calendario del taller.

\section{Arquitectura y Modelado de Datos}
El núcleo del backend se basa en el framework \textbf{NestJS}, estructurado de forma modular, apoyado por \textbf{Prisma ORM} para la conectividad y validación de tipos contra la base de datos relacional (PostgreSQL).

\subsection{Modelos de Prisma (Esquema)}
Las entidades principales desarrolladas son:
\begin{itemize}
    \item \textbf{Usuario:} Maneja atributos de autenticación (email, contraseña encriptada) y de identidad (nombre, teléfono). Implementa un campo \texttt{rol} (\texttt{cliente} o \texttt{mecanico}).
    \item \textbf{Cita:} Representa el agendamiento de un cliente. Se relaciona fuertemente con el modelo Auto y Usuario, controlando atributos como \texttt{fecha\_inicio}, \texttt{descripcion\_problema}, y una robusta \textit{máquina de estados} a través del campo \texttt{estado}.
    \item \textbf{TipoServicio:} Unidades de catálogo de trabajo (ej. Cambio de Aceite, Afinación Mayor) con duraciones paramétricas estandarizadas en días u horas.
    \item \textbf{Bloqueo:} Sistema de exclusividad en calendario que permite a los mecánicos inhabilitar operaciones en ciertas fechas (por vacaciones, remodelaciones, etc.).
\end{itemize}

\section{Mecanismos de Seguridad (Auth \& RBAC)}
Se implementó un estricto flujo de seguridad y aislamiento de datos entre sesiones concurrentes:
\begin{itemize}
    \item \textbf{Autenticación JWT:} Toda petición requiere un token Bearer emitido tras comprobar credenciales (Bcrypt). El ID del usuario y su rol se codifican en el \textit{payload} (\texttt{sub}).
    \item \textbf{Inyección Dinámica:} Se configuró un manejador de ciclo de vida (\texttt{req.user}) para que el ID de quien hace la operación provenga de la firma del servidor, \textbf{bloqueando vulnerabilidades IDOR} donde un usuario podría alterar citas mediante IDs forjados.
    \item \textbf{Controladores de Rol (RBAC):} A través del decorador \texttt{@Roles('mecanico')}, endpoints destructivos o gerenciales quedan blindados ante clientes ordinarios.
\end{itemize}

\section{Flujos y Reglas de Negocio}

\subsection{Máquina de Estados de la Cita}
Se desarrolló una lógica de transiciones estricta para el ciclo de vida del servicio en el taller.
Las transiciones válidas controladas por el sistema son:
\begin{enumerate}
    \item \textbf{pendiente} $\rightarrow$ [aceptada, rechazada, cancelada]
    \item \textbf{aceptada} $\rightarrow$ [en\_curso, cancelada]
    \item \textbf{en\_curso} $\rightarrow$ [completada, cancelada]
    \item \textbf{completada / cancelada / rechazada} $\rightarrow$ \textit{Estado Final (sin transiciones)}
\end{enumerate}

\subsection{Catálogo de Servicios API (Endpoints Relevantes)}

\subsubsection{Citas y Calendario}
\begin{itemize}
    \item \texttt{POST /citas}: Creación de cita por cliente. Vincula el JWT directamente a la entidad creada.
    \item \texttt{GET /citas}: Endpoints polimórficos. Si lo llama un \texttt{cliente}, retorna \textit{solo sus citas}. Si lo llama un \texttt{mecanico}, retorna el \textit{historial completo global}.
    \item \texttt{PATCH /citas/:id}: Orquestador genérico reservado para mecánicos. Modifica \texttt{estado}, \texttt{fecha\_inicio}, \texttt{fecha\_fin} y \texttt{descripcion\_problema} para funciones interactivas como el "drag-and-drop" del dashboard frontend.
    \item \texttt{PATCH /citas/:id/...}: Sub-rutas específicas para cambios ágiles (\texttt{/aceptar}, \texttt{/completar}, \texttt{/cancelar}).
\end{itemize}

\subsubsection{Configuración y Disponibilidad}
\begin{itemize}
    \item \texttt{POST /configuracion/seed}: Automatización de inicialización para cargar el catálogo base (Motor, Cajas, Frenos, etc.).
    \item \textbf{CRUD de Bloqueos:} Conjunto de rutas (\texttt{GET}, \texttt{PATCH}, \texttt{DELETE}) en \texttt{/configuracion/bloqueos} con protección RBAC para bloquear días operativos en el servidor.
\end{itemize}

\section{Testing y Calidad de Integración}
Para asegurar que no hubiera problemas de regresión post-despliegue, se construyó:
\begin{itemize}
    \item Un script interactivo (\texttt{test\_api.py}) que valida mediante heurísticas los códigos HTTP \texttt{200}/\texttt{201} y prueba el flujo integral de un mecánico y un cliente, detectando excepciones (estado 404/500).
    \item Validaciones en el controlador para resolver conflictos de \textit{path routing} estricto en NestJS (Priorizando sub-rutas estáticas sobre dinámicas atrapadas en \texttt{:id}).
\end{itemize}

\section{Infraestructura y Despliegue}
El sistema está paquetizado de principio a fin, habilitando una entrega de integración continua (CI/CD) funcional:
\begin{enumerate}
    \item \textbf{Dockerización:} El entorno empaqueta Node.js junto con los esquemas de Prisma generados en tiempo de construcción.
    \item \textbf{Auto-Migración (\texttt{db push}):} El controlador de entrada (\texttt{CMD}) del \textit{Dockerfile} incrusta el comando \texttt{npx prisma db push}. Cada vez que Render reinicia la instancia tras un nuevo "commit", el esquema refleja e inyecta dinámicamente columnas o tablas nuevas a la base PostgreSQL en la nube, garantizando persistencia y paridad.
    \item \textbf{Hosting:} Desplegado activamente como servicio Web permanente.
\end{enumerate}

\end{document}
